% Generated draft: requires mathematical and source review.
\documentclass[11pt]{article}
\usepackage[margin=25mm]{geometry}
\usepackage{fontspec}
\setmainfont{Latin Modern Roman}
\usepackage{amsmath,amssymb,mathtools,unicode-math}
\setmathfont{Latin Modern Math}
\usepackage{enumitem,needspace,xurl,hyperref}
\setlength{\parindent}{0pt}
\setlength{\parskip}{5pt}
\setlength{\emergencystretch}{3em}
\newcommand{\N}{\mathbb{N}}
\newcommand{\Z}{\mathbb{Z}}
\newcommand{\Q}{\mathbb{Q}}
\newcommand{\R}{\mathbb{R}}
\newcommand{\C}{\mathbb{C}}
\newcommand{\F}{\mathbb{F}}
\newcommand{\A}{\mathbb{A}}
\newcommand{\PP}{\mathbb{P}}
\newcommand{\E}{\mathbb{E}}
\newcommand{\eps}{\varepsilon}
\newcommand{\dd}{\,\mathrm{d}}
\begin{document}
\section*{UnsolvedMath ID: 30005136}
\textbf{Source problem code:} catalogue entry
\subsection{Definitions and mathematical statement}
For each spatial dimension \(n\in\{2,3,4\}\), let \(\mathcal C_n\) denote the family described in the supplied registry as the \(2n\)-parameter family of Lorentzian catenoids. The record does not specify its parametrization, the governing evolution equation, or the initial-data space. Let \(\mathcal X_n\) provisionally denote that unspecified initial-data space, equipped with the intended weighted \(L^2\) topology; this notation does not select a weight, regularity class, or notion of solution.

The supplied mathematical question is: for every \(n\in\{2,3,4\}\), is \(\mathcal C_n\) codimension-one stable under arbitrary nonsymmetric perturbations in the appropriate weighted \(L^2\) topology? Here the intended perturbations are not required to preserve a symmetry. The codimension-one qualification requires a restriction on the admissible perturbation data, but the record does not define that restriction or the ambient space in which its codimension is measured. Thus the question cannot be interpreted as asserting stability for every sufficiently small perturbation without a codimension-one condition.

The source excerpt specifically cautions that translations and Lorentz boosts in the direction of the symmetry axis, as well as scalings, produce perturbations far from the original catenoid at spatial infinity in weighted \(L^2\) topologies. These transformations therefore cannot silently be included as small perturbations in the intended topology. The excerpt also describes \(n=2\) as the analytically most difficult case and says that the methods of its reference [4] extend to higher dimensions; reference [4] is not identified in the supplied record.

A fully quantified stability proposition remains unresolved at the level of formulation: the record supplies neither the weight and data norm nor the evolution, solution class, codimension-one condition, time direction, or stability conclusion. In particular, it does not specify whether stability means remaining close to the family, convergence to a member of the family, or convergence after modulation. No particular choice of these missing ingredients is asserted here.
\subsection{Short English statement}
For each spatial dimension \(2\leq n\leq4\), the supplied question asks whether the \(2n\)-parameter family of Lorentzian catenoids is stable against small perturbations without symmetry assumptions, subject to a codimension-one condition in an appropriate weighted \(L^2\) setting. The precise evolution, topology, condition, and meaning of stability are missing. Supplied dated curation labels the problem partially solved, reporting nonsymmetric nonlinear asymptotic stability for \(n\geq3\) while identifying the two-dimensional nonradial case and stability of general stationary solutions as open; these claims are not independently verified here.
\subsection{Sources}
\begin{itemize}
\item Local UnsolvedMath registry, supplied complete record for permanent ID \(30005136\), title \emph{Codimension-One Stability of Lorentzian Catenoids}, exact source problem code \texttt{catalogue entry}. No individual proposer is supplied. The proposed year is recorded as 2022.
\item Supplied report citation: \emph{Deterministic Dynamics and Randomness in PDE (2023)}, Oberwolfach Reports, pp.\ 1431--1499, DOI \texttt{10.4171/owr/2022/26}. \url{https://doi.org/10.4171/owr/2022/26}. The background describes Oberwolfach Reports as documenting mathematical workshops and their open-problem discussions. The registry labels its statement \texttt{corrected\_verified} and reports checking it against the official report; that verification is supplied provenance, not an independent check performed here.
\item Supplied literature source: Luehrmann--Oh--Shahshahani, catenoid stability. \url{https://doi.org/10.1007/s00222-025-01330-3}.
\item Supplied literature source: Tang, four-dimensional catenoid. \url{https://arxiv.org/abs/2411.08366}.
\item Supplied literature source: Oberwolfach report. \url{https://doi.org/10.4171/owr/2022/26}.
\item Literature assessment dated exactly \texttt{2026-08-22}: the supplied curation labels the problem partially solved and reports nonlinear asymptotic stability without symmetry for spatial dimensions \(n\geq3\), with the two-dimensional nonradial case and stability of general stationary solutions remaining open. This is supplied dated curation derived from the OWR material, not independently verified literature status. Its time-sensitive status and cited claims require independent verification.
\end{itemize}
\end{document}
